Reasoning large language models can be copied by distilling chain-of-thought traces
into cheaper students, but provenance tests must separate descent from ordinary task
competence. We study a forensic audit in which the suspect is observed only through
output traces, candidate teachers are available for authorized likelihood scoring,
and the suspect base is known or declared. The common base-relative likelihood
surplus is confounded by candidate capability: in our GSM8K audit it false-alarms
on a human-reference student (empirical false-positive rate, FPR, $0.81$) and misses
same-family distillation (area under the receiver-operating characteristic curve,
AUROC, $0.13$). Capability-matched likelihood (CML) instead compares suspect traces
with a same-base non-distilled reference under the same candidate teacher. On
three-seed Qwen2.5-7B-Instruct GSM8K students distilled from three teachers, CML
reaches trace-level AUROC $0.997$, true positive rate at $1\%$ FPR $0.991$, and
empirical FPR $0.2\%$; it also recovers same-family distillation (AUROC $0.998$). A
single-seed MATH stress test shows the same pattern (AUROC $0.993$). Multi-Task
Calibration Reference (MTCR) improves sibling-teacher attribution from
$0.672$/$0.658$ to $0.894$/$0.875$. Aggregate stress summaries further test adaptive
trace transformations, cross-base and cross-task transfer, open-set decoys, and
student-level confidence intervals: CML+MTCR retains AUROC at least $0.930$ across
the reported adaptive sweeps, the cross-base/task matrix has minimum AUROC $0.912$,
and open-set false attribution stays at or below $2.0\%$.
